\RequirePackage{pdfmanagement}
\SetKeys[document/metadata]{pdfversion=1.7, pdfstandard=A-2u, lang=en-GB}
\input{preamble/pdfA-compliance/glyphtounicode}
\documentclass[11pt,a4paper]{report}

\newcommand*{\PaperTitle}{Absolute concentration robustness with multistability in bimolecular networks: the exact minimal sizes}
\newcommand*{\PaperTitlePlain}{Absolute concentration robustness with multistability in bimolecular networks: the exact minimal sizes}
\newcommand*{\PaperShortTitle}{ACR with multistability}
\newcommand*{\PaperSubtitle}{chemistry-open-problems \textperiodcentered{} 31 October 2026 \textperiodcentered{} merged manuscript, verification pending}
\newcommand*{\PaperCredits}{Solvers: Claude Opus 5.5 (effort xhigh; dimension two, species, complexes), GPT-6-Astra (effort xhigh; five reactant complexes) \textperiodcentered{} Merge: Claude Opus 5.5 (effort not recorded) \textperiodcentered{} Verifiers of the two parts: Claude Opus 5.5 (effort xhigh; effort not recorded) \textperiodcentered{} Verification of the merged text: pending}
\newcommand*{\PaperKeywords}{absolute concentration robustness, multistability, reaction networks, mass-action kinetics, bimolecular networks, deficiency, exact enumeration, toric steady states}
\input{chemop-preamble}
\addbibresource{references.bib}

\providecommand{\R}{\mathbb{R}}
\newcommand{\Rpos}{\mathbb{R}_{>0}}
\newcommand{\Rnn}{\mathbb{R}_{\ge 0}}
\newcommand{\spn}{\operatorname{span}}
\newcommand{\supp}{\operatorname{supp}}
\newcommand{\one}{\mathbf{1}}
\newcommand{\Vp}{V^{+}}
\providecommand{\rank}{\operatorname{rank}}
\providecommand{\im}{\operatorname{im}}
\providecommand{\diag}{\operatorname{diag}}
\providecommand{\sgn}{\operatorname{sgn}}

\begin{document}
\makepapertitle

\begin{paperabstract}
Kaihnsa, Nguyen and Shiu proved that an at-most-bimolecular reaction network in which absolute concentration
robustness (ACR) and nondegenerate multistationarity coexist has at least 3 species, 3 reactant complexes,
5 complexes and a two-dimensional stoichiometric subspace, all four bounds being attained. They conjectured that
each of these four lower bounds becomes strictly larger when multistationarity is replaced by multistability. We
determine the four exact minima for ACR with multistability: dimension $3$, $3$ species, $5$ reactant complexes and
$6$ complexes. The conjecture therefore holds for three of the four quantities and fails for the number of species.
Two theorems carry the lower bounds. A two-dimensional bimolecular mass-action system with ACR has at most one
exponentially stable positive steady state in each stoichiometric compatibility class; the proof combines a normal
form of the vector field on the ACR hyperplane, a connectedness property of the positive steady-state set that holds
precisely because products are at most bimolecular, and the global ACR condition. A rank-three system with four
reactant complexes and ACR has at most two nondegenerate positive steady states per class, of opposite determinant
signs; an exact finite classification reduces this to nineteen scalar normal forms. Six complexes follow by a
kinetic-subspace argument and the deficiency-one theorem. Two explicit networks attain the bounds. The six-reaction
network $\{Z\to Y+Z,\ 2Z\to Z,\ Y\to Y+Z,\ Y+Z\to 2X,\ X+Z\to 2Z,\ 2X\to Y\}$ has ACR in $X$ for all rate constants
and, for suitable ones, two exponentially stable positive steady states, with $3$ species, $6$ complexes and
dimension $3$. A four-species, six-reaction network with integer rates has $5$ reactant complexes. Stability and ACR
are certified in exact arithmetic. We do not know whether one network attains all four minima.
\end{paperabstract}

\section*{The question}

Kaihnsa, Nguyen and Shiu~\autocite{KNS2024} study the minimal size of reaction networks in which ACR and
multistationarity coexist. Their main result (Theorem~1.2 there) states that an at-most-bimolecular network~$G$
admitting a vector of rate constants for which the mass-action system has ACR and is nondegenerately
multistationary has at least 3 species, at least 3 reactant complexes, at least 5 complexes and dimension at least
2; the network $\{X_1+X_2\to 2X_3,\ X_3\to X_1,\ 2X_3\to 2X_2\}$ attains all four bounds. Section~6 of the paper
ends with the following conjecture:

\begin{openquote}
``In the present work, our interest in multistationarity comes from the fact that it is a necessary condition for
multistability. Another interesting direction, therefore, is to investigate conditions for coexistence of ACR and
multistability, rather than multistationarity. The "minimal" networks in the current work admit only two positive
steady states and are not multistable. Hence, we conjecture that the lower bounds (on dimension and the numbers of
species, reactant complexes, and all complexes) for the coexistence of ACR and multistability are strictly larger
than the bounds proven here for ACR and multistationarity.''
\end{openquote}

\begin{theorem}[main result]\label{thm:main}
Let $G$ be an at-most-bimolecular reaction network and suppose that for some vector of positive rate constants
$\kappa$ the mass-action system $(G,\kappa)$ has ACR and is multistable. Then
\begin{enumerate}
\item the stoichiometric subspace of $G$ has dimension at least $3$;
\item $G$ has at least $3$ species;
\item $G$ has at least $5$ reactant complexes;
\item $G$ has at least $6$ complexes.
\end{enumerate}
All four bounds are attained: the network $\mathcal N_1$ of \zcref{thm:three-species}, which has ACR for every
vector of rate constants, attains those on dimension, species and complexes, and the network $\mathcal N_2$ of
\zcref{sec:five} attains the one on reactant complexes. Hence the conjecture holds for dimension, reactant complexes
and complexes, and fails for the number of species.
\end{theorem}

\begin{table}[ht]
\centering
\begin{tabular}{lccc}
\toprule
& ACR + nondegenerate & ACR + multistability & attained\\
& multistationarity \autocite{KNS2024} & exact minimum & by\\
\midrule
dimension & $2$ & $3$ & $\mathcal N_1$\\
species & $3$ & $3$ & $\mathcal N_1$\\
reactant complexes & $3$ & $5$ & $\mathcal N_2$\\
complexes & $5$ & $6$ & $\mathcal N_1$\\
\bottomrule
\end{tabular}
\caption{Minimal sizes of at-most-bimolecular networks. All bounds of the first column are attained by the network
$\{X_1+X_2\to2X_3,\ X_3\to X_1,\ 2X_3\to2X_2\}$ of \autocite{KNS2024}.}
\label{tab:summary}
\end{table}

Part~1 follows from \zcref{thm:dim-two} (\zcref{cor:dimension}), and part~2 from \autocite[Theorem~4.1]{KNS2024},
because a multistable system is nondegenerately multistationary. Part~3 combines \zcref{thm:dim-two} with
\zcref{thm:rank-three} (\zcref{cor:five}). Part~4 is \zcref{cor:complexes}. The lower and upper bounds concern
arbitrary numbers of species and reactions, and no claim of simultaneous minimization is made.

This text merges two reports of this repository, each checked by a fresh-context verifier: one on the conjecture
(sections on dimension two, complexes and three species), one on the exact number of reactant complexes. Their
original versions are kept in the folder \texttt{archive/}. The merged text has not yet been re-verified.

\section{Setting}\label{sec:setting}

A reaction network $G$ on species $X_1,\dots,X_n$ is a finite set of reactions $y\to y'$ between complexes
$y,y'\in\mathbb{Z}_{\ge0}^n$, $y\neq y'$. It is \emph{at most bimolecular} (bimolecular for short) if every
complex, reactant or product, satisfies $|y|:=y_1+\dots+y_n\le 2$. Given positive rate constants
$\kappa=(\kappa_{y\to y'})$, the mass-action system is
\begin{equation}\label{eq:ma}
\dot x=f(x):=\sum_{y\to y'}\kappa_{y\to y'}\,x^{y}\,(y'-y)=\sum_{y}x^{y}\,v_{y},\qquad
v_{y}:=\sum_{y'\colon y\to y'}\kappa_{y\to y'}(y'-y),
\end{equation}
the second sum running over reactant complexes. The stoichiometric subspace is
$S=\spn\{y'-y\}$, $s=\dim S$ is the dimension of $G$, and the stoichiometric compatibility classes are the sets
$(x^0+S)\cap\Rnn^n$. A positive steady state $x^*$ is \emph{nondegenerate} if the restriction $Df(x^*)|_S\colon S\to
S$ is invertible, and \emph{exponentially stable} if all eigenvalues of $Df(x^*)|_S$ have negative real part; the
latter implies the former.

\begin{definition}\label{def:acr-ms}
Following \autocite[Definition~2.6]{KNS2024}, $(G,\kappa)$ has \emph{ACR in $X_i$} if it has a positive steady state
and $x_i$ takes the same value $c>0$ at all positive steady states (in all compatibility classes). We call
$(G,\kappa)$ \emph{multistable} if some compatibility class contains at least two exponentially stable positive
steady states. Both properties refer to the same $\kappa$.
\end{definition}

This is the usual notion of stability in the literature on small multistable networks
\autocite[§2.2]{JTZ2025}: a steady state is stable if it is nondegenerate and the nonzero eigenvalues of the Jacobian
have negative real parts, which is equivalent to our definition. As in \autocite[Theorem~1.2]{KNS2024}, the
question is existential: a network qualifies if one vector $\kappa$ gives both properties.

Exponentially stable steady states are nondegenerate, so a multistable system is nondegenerately multistationary and
every lower bound of \autocite{KNS2024} applies to ACR with multistability.

The zero complex is allowed and counts as a complex, also as a reactant complex: Example~4.13
of~\autocite{KNS2024} counts $0\to X$ among its four reactant complexes. The result on reactant complexes depends on
this convention. Arbitrarily many reactions may leave each reactant complex.

\begin{lemma}[coefficient cones]\label{lem:cones}
Let $G$ be bimolecular and let $v_y$ be the coefficient vector of the monomial $x^y$ in~\eqref{eq:ma}. Then
$(v_y)_i\ge 0$ for every $i\notin\supp(y)$, and $\sum_i (v_y)_i\le 0$ whenever $|y|=2$.
\end{lemma}

\begin{proof}
Each $y'-y$ with $y\to y'$ has $(y'-y)_i=y'_i\ge 0$ for $i\notin\supp(y)$, and $\sum_i(y'-y)_i=|y'|-|y|\le 0$ if
$|y|=2$, because $|y'|\le 2$. Both properties pass to nonnegative combinations.
\end{proof}

The first property is the familiar characterization of mass-action vector fields~\autocite[Lemma~2.3]{KNS2024}; the
second one is where bimolecularity of the \emph{products} enters, and it is essential below
(\zcref{rem:products}).

\section{Dimension two}\label{sec:dim-two}

Throughout this section $G$ is bimolecular with $\dim S=2$, $\kappa$ is fixed, $(G,\kappa)$ has ACR in $X_1$ with
value $c>0$, $f$ is the vector field~\eqref{eq:ma} and $H:=\{x\in\R^n\colon x_1=c\}$. All positive steady states lie
in $H$. For $x\in H$ we write $x=(c,x')$ with $x'\in\R^{n-1}$.

\begin{lemma}\label{lem:not-conserved}
If $(G,\kappa)$ has a nondegenerate positive steady state, then $e_1\notin S^\perp$.
\end{lemma}

\begin{proof}
Let $x^*$ be a nondegenerate positive steady state and let $W$ be a matrix whose $n-2$ rows form a basis of
$S^\perp$. The map $F(x)=(f(x),Wx)\in S\times\R^{n-2}$ has invertible derivative at $x^*$: if $Df(x^*)u=0$ and
$Wu=0$ then $u\in S$, hence $u=0$. By the inverse function theorem, for every $T$ near $Wx^*$ there is a positive
steady state $x(T)$ near $x^*$ with $Wx(T)=T$. If $e_1\in S^\perp$ we may take $e_1^{\mathsf T}$ as a row of $W$; the
corresponding coordinate of $T$ equals $x_1(T)$, which can then be made different from $c$, contradicting ACR.
\end{proof}

Under the hypothesis of the lemma, $S\cap e_1^\perp$ is one-dimensional; we fix a vector $\ell$ spanning it. Every
compatibility class meets $H$ in a segment parallel to~$\ell$.

\begin{lemma}[the field on the ACR hyperplane]\label{lem:field-on-H}
Suppose some compatibility class contains two distinct nondegenerate positive steady states $p$ and $q$. Then there
are $\alpha\in S\setminus\{0\}$ and a polynomial $h$ of degree at most $2$ such that $f(x)=h(x)\,\alpha$ for all
$x\in H$. Moreover $q=p+\tau\ell$ for some $\tau\neq0$, and $h_L(t):=h(p+t\ell)$ has degree exactly $2$ with simple
zeros $t=0$ and $t=\tau$.
\end{lemma}

\begin{proof}
Choose a basis $\sigma_1,\sigma_2$ of $S$ and write $f=g_1\sigma_1+g_2\sigma_2$ with polynomials $g_1,g_2$ of degree
at most $2$. Since $Df(p)$ maps onto $S$, the set $f^{-1}(0)$ is, near $p$, a smooth submanifold of dimension $n-2$,
and by ACR its positive points lie in~$H$. Thus the restrictions $g_1|_H$ and $g_2|_H$, viewed as polynomials on
$H\cong\R^{n-1}$, vanish on a smooth hypersurface $M$ of $H$.

We claim that they have a nonconstant common factor (or one of them vanishes identically, in which case the lemma
holds with $\alpha$ the other basis vector). By Baire's theorem, an irreducible factor $F_1$ of $g_1|_H$ vanishes on a
nonempty open subset $M_1$ of $M$, and an irreducible factor $F_2$ of $g_2|_H$ vanishes on a nonempty open subset
$M_2\subseteq M_1$. If $F_1$ and $F_2$ were not proportional, they would be coprime, their common complex zero set
would have dimension at most $n-3$, and so would its real points; but it contains the $(n-2)$-dimensional manifold
$M_2$.

If the common factor $F$ has degree $2$, then $g_i|_H=c_iF$ with constants $c_i$, and $f|_H=F\,(c_1\sigma_1+c_2\sigma_2)$.
If $F$ has degree $1$, then $g_i|_H=F\,m_i$ with $\deg m_i\le 1$. If $m_1,m_2$ are proportional we conclude as before.
Otherwise $\{x\in H\colon m_1=m_2=0\}$ is empty or an affine subspace of dimension at most $n-3$. A nondegenerate
positive steady state $y$ with $F(y)\neq 0$ would have a neighbourhood in $f^{-1}(0)$, of dimension $n-2$, contained
in that set; so every nondegenerate positive steady state satisfies $F=0$. The steady states of the class of $p$ lie
on the segment $(p+\R\ell)\cap\Rpos^n$, on which $F$ is affine; vanishing at $p\neq q$, it vanishes on the whole
segment, so $p$ is not isolated in its class, which contradicts nondegeneracy (nondegenerate steady states are
isolated in their class by the inverse function theorem).

Finally, $p,q\in H$ lie in the same class, so $q-p\in S\cap e_1^\perp=\R\ell$. The polynomial $h_L$ vanishes at $0$
and $\tau$; if its degree were at most one it would vanish identically and $p$ would again not be isolated. Hence
$\deg h_L=2$ and both zeros are simple. Likewise $f|_H\not\equiv0$ (otherwise $p$ would not be isolated), so the
vector $\alpha$ obtained above is nonzero.
\end{proof}

From now on we assume that a class contains two distinct nondegenerate positive steady states $p,q$, and we use
$\alpha$, $h$, $\ell$, $\tau$ and $h_L$ from \zcref{lem:field-on-H}.

\begin{lemma}\label{lem:alpha1-zero}
If $\alpha_1=0$, then $Df|_S$ has a positive real eigenvalue at $p$ or at $q$.
\end{lemma}

\begin{proof}
Then $\alpha\in S\cap e_1^\perp$, so $\alpha=\gamma\ell$ with $\gamma\neq0$. For $x_t:=p+t\ell\in H$ we have
$f(x_t)=\gamma h_L(t)\ell$, hence $Df(x_t)\ell=\gamma h_L'(t)\ell$: $\ell$ is an eigenvector with eigenvalue
$\gamma h_L'(t)$. Since $h_L'(0)$ and $h_L'(\tau)$ are nonzero with opposite signs, this eigenvalue is positive at $p$
or at $q$.
\end{proof}

\paragraph{Normal form.}
If $\alpha_1\neq0$ we normalize $\alpha_1=1$. Since $\{\alpha,\ell\}$ is a basis of $S$ and $\ell_1=0$, we can write
uniquely
\begin{equation}\label{eq:normal-form}
f=\psi\,\alpha+\varphi\,\ell,\qquad \psi=f_1 .
\end{equation}
On $H$ we have $f=h\alpha$, so $\varphi$ vanishes on $H$ and $\psi|_H=h$. As $\deg\varphi\le2$, there is an affine
function $P$ with $\varphi=(x_1-c)\,P$. The positive steady states are the positive solutions of $\psi=0$ and
$(x_1-c)P=0$.

\begin{lemma}[determinant on the ACR hyperplane]\label{lem:det}
If $\alpha_1=1$, then $\det\bigl(Df(x)|_S\bigr)=-P(x)\,\bigl(\nabla\psi(x)\cdot\ell\bigr)$ for every $x\in H$.
\end{lemma}

\begin{proof}
On $H$, $Df=\alpha\,\nabla\psi^{\mathsf T}+P\,\ell\,e_1^{\mathsf T}$. Hence $Df\,\alpha=(\nabla\psi\cdot\alpha)\alpha+P\ell$
and $Df\,\ell=(\nabla\psi\cdot\ell)\alpha$. In the basis $(\alpha,\ell)$ the matrix of $Df|_S$ is
$\begin{psmallmatrix}\nabla\psi\cdot\alpha&\nabla\psi\cdot\ell\\ P&0\end{psmallmatrix}$, with determinant
$-P\,(\nabla\psi\cdot\ell)$.
\end{proof}

\begin{corollary}\label{cor:P-changes-sign}
If $\alpha_1=1$ and $p,q$ are both exponentially stable, then $P(p)\,P(q)<0$.
\end{corollary}

\begin{proof}
A $2\times2$ matrix whose eigenvalues have negative real parts has positive determinant. Here
$\nabla\psi(p)\cdot\ell=h_L'(0)$ and $\nabla\psi(q)\cdot\ell=h_L'(\tau)$ are nonzero with opposite signs, so
\zcref{lem:det} forces $P(p)$ and $P(q)$ to have opposite signs.
\end{proof}

The remaining steps show that $P$ cannot change sign along the positive steady states. Write
$\psi_H(x'):=\psi(c,x')$ and
\[
\Vp:=\{x'\in\Rpos^{n-1}\colon \psi_H(x')=0\},
\]
so that the positive steady states are exactly the points $(c,x')$ with $x'\in\Vp$.

\begin{lemma}[bimolecularity makes $\Vp$ connected]\label{lem:connected}
Assume $\alpha_1=1$. Then $\Vp$ is connected, and $\nabla\psi_H\neq0$ at every point of $\Vp$.
\end{lemma}

\begin{proof}
For $2\le j\le k\le n$ let $a_{jk}$ be the coefficient of $x_jx_k$ in $\psi=f_1$. Monomials containing $x_1$ become
of degree at most one in $x'$ on $H$, and $(x_1-c)P$ contains no monomial $x_jx_k$ with $j,k\ge2$; hence the
quadratic part of $\psi_H$ is $Q(x')=\sum_{j\le k}a_{jk}x_jx_k$ and the coefficient vector of $x_jx_k$ in $f$ is
$a_{jk}\alpha$. By \zcref{lem:cones} applied to $y=X_j+X_k$: $a_{jk}\ge0$ (coordinate $1$), and if $a_{jk}>0$ then
$\alpha_i\ge0$ for all $i\notin\{1,j,k\}$ and $1+\sum_{i\ge2}\alpha_i\le 0$, so that
$\sum_{i\in\{j,k\}}\alpha_i\le -1-\sum_{i\notin\{1,j,k\}}\alpha_i\le -1$ (a single term when $j=k$). Let
$N:=\{i\ge2\colon\alpha_i<0\}$. Every monomial
$x_jx_k$ with $a_{jk}>0$ therefore satisfies $N\subseteq\{j,k\}$ and $N\cap\{j,k\}\neq\emptyset$. Since
$h_L(t)=\psi_H(p'+t\ell')$ has degree $2$, $Q(\ell')\ne0$, so $Q\neq0$ and $1\le|N|\le2$.

Second, for $i\in N$ the component $f_i$ restricted to $H$ equals $\alpha_i\psi_H$; its monomials not containing
$x_i$ have nonnegative coefficients (setting $x_1=c>0$ preserves this property of $f_i$), so the monomials of
$\psi_H$ not containing $x_i$ have nonpositive coefficients. Two cases remain.

\emph{Case $N=\{a,b\}$.} Then $Q=\beta x_ax_b$ with $\beta>0$, every other monomial of $\psi_H$ lacks $x_a$ or $x_b$,
and $\psi_H=\beta x_ax_b-G_0-G_ax_a-G_bx_b-\sum_{k\notin\{a,b\}}G_kx_k$ with all $G\ge0$. Writing
$F:=G_0+G_ax_a+\sum_{k\notin\{a,b\}}G_kx_k$, the equation $\psi_H=0$ reads $x_b(\beta x_a-G_b)=F$. On the open
orthant either $F>0$, and then $\Vp$ is the graph $x_b=F/(\beta x_a-G_b)$ over the convex set
$\{\beta x_a>G_b\}\cap\Rpos^{n-2}$, with $\partial\psi_H/\partial x_b=\beta x_a-G_b>0$; or $F\equiv0$, and then
$\Vp=\{x_a=G_b/\beta\}\cap\Rpos^{n-1}$ is convex, with $\partial\psi_H/\partial x_a=\beta x_b>0$.

\emph{Case $N=\{a\}$.} Then $Q=x_a\sum_{k\ge2}\beta_kx_k$ with all $\beta_k\ge0$, not all zero, and
$\psi_H=\beta_ax_a^2+x_a\bigl(K+R(x'')\bigr)-G(x'')$, where $x''$ collects the coordinates other than $x_1,x_a$,
$R(x'')=\sum_{k\neq a}\beta_kx_k$, $K\in\R$, and $G$ is affine with nonnegative coefficients. On the open orthant
either $G>0$ or $G\equiv0$. If $G>0$ and $\beta_a>0$, the quadratic in $x_a$ has exactly one positive root, which is
simple; if $G>0$ and $\beta_a=0$, the positive root $x_a=G/(K+R)$ exists exactly on the convex set $\{K+R>0\}$, where
$\partial\psi_H/\partial x_a=K+R>0$. In both cases $\Vp$ is the graph of a smooth function of $x''$ over a convex set.
If $G\equiv0$, then $\Vp=\{\beta_ax_a+K+R(x'')=0\}\cap\Rpos^{n-1}$ is convex, and on it
$\nabla\psi_H=x_a\nabla(\beta_ax_a+K+R)\neq0$, because $\beta_a$ or some $\beta_k$ with $k\ne a$ is positive.
\end{proof}

\begin{lemma}[global ACR]\label{lem:global}
Assume $\alpha_1=1$ and $P(p)P(q)<0$, and write $P_H(x'):=P(c,x')$. Then
$Z:=\{x'\in\Rpos^{n-1}\colon P_H(x')=0\}$ is nonempty, and there is $\varepsilon\in\{1,-1\}$ with
$\varepsilon\,\psi_H\ge0$ on $Z$.
\end{lemma}

\begin{proof}
$P_H$ takes both signs on the convex set $\Rpos^{n-1}$, so $Z\neq\emptyset$; in particular $P$ depends on some
$x_k$ with $k\ge2$, and the hyperplane $\Pi:=\{P=0\}$ of $\R^n$ is transverse to $H$. At a positive point of $\Pi$
we have $f=\psi\alpha$, so by ACR $\psi$ has no zero on the convex sets
$\Pi^{\pm}:=\Pi\cap\Rpos^n\cap\{\pm(x_1-c)>0\}$, and hence has a constant sign $\varepsilon_\pm$ on each. Both sets
accumulate at every point of $\{c\}\times Z$ (move inside $\Pi$ along $e_1-(\pi_1/\pi_k)e_k$, where $\pi_j$ are the
coefficients of $P$ and $\pi_k\neq0$), so $\varepsilon_+\psi_H\ge0$ and $\varepsilon_-\psi_H\ge0$ on $Z$. If
$\varepsilon_+\neq\varepsilon_-$, then $\psi_H$ vanishes on $Z$, a nonempty open subset of the hyperplane
$\{P_H=0\}$, so $P_H$ divides $\psi_H$: $\psi_H=P_H\,A$ with $A$ affine. Since $P_H(p'),P_H(q')\neq0$, $A$ vanishes at
$p'$ and $q'$, hence on the line through them, and $h_L\equiv0$, contradicting \zcref{lem:field-on-H}. So
$\varepsilon:=\varepsilon_+=\varepsilon_-$ works.
\end{proof}

\begin{lemma}[no transversal sign change]\label{lem:tangency}
In the situation of \zcref{lem:global}, let $z\in\Vp\cap Z$. Then $\nabla\psi_H(z)=\mu\,\nabla P_H$ for some
$\mu\neq0$, and $z$ has a neighbourhood $U$ in $\Vp$ on which $-\operatorname{sign}(\mu)\,\varepsilon\,P_H\ge0$.
\end{lemma}

\begin{proof}
Let $E:=\{P_H=0\}$, with direction space $T$. The function $\varepsilon\psi_H$ is nonnegative on the relatively open
set $Z\subseteq E$ and vanishes at $z$, so $\nabla\psi_H(z)\perp T$, that is $\nabla\psi_H(z)=\mu\nabla P_H$; and
$\mu\neq0$ by \zcref{lem:connected}. Let $Q$ be the quadratic part of $\psi_H$ and $B$ its polar form. For $v\in T$,
$\psi_H(z+v)=Q(v)$, so $\varepsilon Q\ge0$ on $T$. Fix $\nu$ with $\nabla P_H\cdot\nu=1$ and write $w=v+s\nu$ with
$v\in T$, $s=P_H(z+w)$. If $z+w\in\Vp$,
\[
0=\psi_H(z+w)=\mu s+Q(v)+2sB(v,\nu)+s^2Q(\nu),\quad\text{i.e.}\quad s\,\bigl(\mu+2B(v,\nu)+sQ(\nu)\bigr)=-Q(v).
\]
For $w$ small the bracket has the sign of $\mu$, so $s$ is zero or has the sign of $-\mu\varepsilon$.
\end{proof}

\begin{theorem}[dimension two]\label{thm:dim-two}
Let $G$ be an at-most-bimolecular network of dimension $2$ and let $\kappa$ be such that $(G,\kappa)$ has ACR. Then
every stoichiometric compatibility class contains at most one exponentially stable positive steady state. In
particular $(G,\kappa)$ is not multistable.
\end{theorem}

\begin{proof}
Suppose $p\neq q$ are exponentially stable positive steady states in one class; they are nondegenerate. By
\zcref{lem:not-conserved,lem:field-on-H} we have the objects $\alpha,\ell,h$; by \zcref{lem:alpha1-zero},
$\alpha_1\neq0$, and we normalize $\alpha_1=1$. By \zcref{cor:P-changes-sign}, $P(p)P(q)<0$. Define
$\theta\colon\Vp\to\{1,-1\}$ by $\theta(x')=\operatorname{sign}P_H(x')$ if $P_H(x')\neq0$ and
$\theta(x')=-\operatorname{sign}(\mu(x'))\,\varepsilon$ if $x'\in Z$, where $\mu(x')=\nabla\psi_H(x')\cdot\nabla
P_H/|\nabla P_H|^2$ is continuous and nonzero on $\Vp\cap Z$ (\zcref{lem:tangency}). Then $\theta$ is locally
constant: near a point with $P_H\neq0$ by continuity, and near $z\in\Vp\cap Z$ by \zcref{lem:tangency} applied on a
neighbourhood where $\mu$ keeps its sign. As $\Vp$ is connected (\zcref{lem:connected}), $\theta$ is constant, so
$P(p)$ and $P(q)$ have the same sign, a contradiction.
\end{proof}

\begin{remark}[where the bimolecular products matter]\label{rem:products}
The bound $|y'|\le2$ on product complexes is used only through the second property in \zcref{lem:cones}, inside
\zcref{lem:connected}. Without it, $\psi_H$ may contain the terms $x_2^2-3x_2+2$, and $\Vp$ may consist of two
parallel sheets on which $P$ has opposite signs; the argument then breaks down. In the planar model
$\dot u=t^2-1-u/c$, $\dot t=-ut$ (with $u=x_1-c$), the two stable nodes $t=\pm1$ of the ACR line $u=0$ are separated
by the invariant line $t=0$, which ends at a saddle on the face $x_1=0$. We have not found a mass-action realization
of this mechanism with bimolecular reactant complexes (linear programs with the constraints of \zcref{lem:cones}
minus the product bound are feasible locally, but the global ACR condition failed on every sample), so we do not
know whether \zcref{thm:dim-two} extends to networks with larger product complexes.
\end{remark}

\section{Dimension and number of complexes}\label{sec:consequences}

\begin{corollary}[dimension]\label{cor:dimension}
If a bimolecular network $G$ admits $\kappa$ with ACR and multistability, then $\dim S\ge3$.
\end{corollary}

\begin{proof}
One-dimensional networks never have ACR with nondegenerate multistationarity \autocite[Proposition~3.10]{KNS2024},
dimension~$2$ is excluded by \zcref{thm:dim-two}, and dimension~$0$ has singleton classes.
\end{proof}

For the number of complexes we use the classical objects of chemical reaction network theory: $m$ complexes, $\ell$
linkage classes, $t$ terminal strong linkage classes, deficiency $\delta=m-\ell-s$, and the factorization
$f(x)=Y A_\kappa \Psi(x)$, where $Y$ is the $n\times m$ matrix of complexes, $\Psi(x)=(x^{y})_y$ and $A_\kappa$ is
the Laplacian matrix of the weighted reaction graph \autocite{Feinberg2019}. The kinetic subspace
$K_\kappa:=Y(\operatorname{Im}A_\kappa)\subseteq S$ contains every value of $f$ and every column of $Df$.

\begin{lemma}[kinetic subspace]\label{lem:kinetic}
If $(G,\kappa)$ has a nondegenerate positive steady state and $G$ has a nonterminal complex, then $\delta\ge
t-\ell+1$.
\end{lemma}

\begin{proof}
At a nondegenerate positive steady state $x^*$ the image of $Df(x^*)|_S$ is $S$, so $K_\kappa=S$. The kernel of
$A_\kappa$ has dimension $t$ and consists of vectors supported on terminal complexes, so
$\dim\operatorname{Im}A_\kappa=m-t$, and $A_\kappa\Psi(x^*)\neq0$ because $\Psi(x^*)$ is positive on a nonterminal
complex. As $YA_\kappa\Psi(x^*)=f(x^*)=0$, the map $Y$ has a nonzero kernel on $\operatorname{Im}A_\kappa$. Hence
$s=\dim K_\kappa\le m-t-1$, that is $\delta=m-\ell-s\ge t-\ell+1$.
\end{proof}

\begin{corollary}[complexes]\label{cor:complexes}
If a bimolecular network $G$ admits $\kappa$ with ACR and multistability, then $G$ has at least $6$ complexes.
\end{corollary}

\begin{proof}
By \zcref{cor:dimension}, $s\ge3$. A network of deficiency zero is not multistationary \autocite[Lemma~2.4]{KNS2024},
so $\delta\ge1$ and $m=\ell+s+\delta\ge5$, with equality only if $\ell=1$, $s=3$, $\delta=1$. Suppose $m=5$. If $G$
has a nonterminal complex, \zcref{lem:kinetic} gives $t\le\ell+\delta-1=1$; otherwise all complexes are terminal and
$t=\ell=1$. So $G$ has one linkage class, of deficiency one, with a single terminal strong linkage class, and the
deficiency-one theorem \autocite{Feinberg1995} gives at most one positive steady state in each compatibility class.
This contradicts multistability.
\end{proof}

\begin{remark}[on the deficiency-one theorem]\label{rem:def-one}
The uniqueness part of the deficiency-one theorem holds without weak reversibility, provided each linkage class
contains exactly one terminal strong linkage class, $\delta_\theta\le1$ for each linkage class and
$\sum_\theta\delta_\theta=\delta$ \autocite{Feinberg1995}. \textcite[Lemma~2.5]{KNS2024} state it for weakly reversible
networks only, but their proof of the bound of $5$ complexes (Theorem~4.9(2) there) applies it to networks that need
not be weakly reversible; \zcref{lem:kinetic} supplies the missing hypothesis $t=1$ there, so their bound stands. The
same argument shows that two nondegenerate positive steady states in one class force $m\ge s+3$ for networks of any
molecularity: $\delta\ge1$ always, and $\delta\ge2$ when $\ell=1$.
\end{remark}

\section{Three species and six complexes suffice}\label{sec:three}

\begin{theorem}[three species, six complexes]\label{thm:three-species}
The bimolecular network on species $X,Y,Z$
\[
\begin{gathered}
\mathcal N_1:\quad Z\xrightarrow{k_1}Y+Z,\quad 2Z\xrightarrow{k_2}Z,\quad Y\xrightarrow{k_3}Y+Z,\\
Y+Z\xrightarrow{k_4}2X,\quad X+Z\xrightarrow{k_5}2Z,\quad 2X\xrightarrow{k_6}Y
\end{gathered}
\]
has $3$ species, $6$ reactions, $6$ reactant complexes, $6$ complexes and dimension $3$. For every vector of rate
constants it has ACR in $X$, with value $2k_1/k_5$. For $(k_1,\dots,k_6)=(5,1,1,4,10,7)$ its positive steady states
are exactly
\[
P_1=\bigl(1,\tfrac{19}{4},\tfrac12\bigr),\qquad P_2=(1,3,1),\qquad P_3=\bigl(1,\tfrac74,\tfrac72\bigr),
\]
and $P_1,P_3$ are exponentially stable while $P_2$ is a saddle. The network being full-dimensional, all three lie in
the unique compatibility class $\Rnn^3$. Multistability and ACR persist for all rate constants in a neighbourhood of
this vector.
\end{theorem}

\begin{proof}
The mass-action system is
\begin{align*}
\dot x&=2k_4yz-k_5xz-2k_6x^2,\\
\dot y&=k_1z-k_4yz+k_6x^2,\\
\dot z&=-k_2z^2+k_3y-k_4yz+k_5xz .
\end{align*}
The reaction vectors $(0,1,0)$, $(0,0,-1)$, $(2,-1,-1)$ and $(-1,0,1)$ span $\R^3$. The weight $w=(1,2,0)$ gives
$\dot x+2\dot y=z\,(2k_1-k_5x)$, so every positive steady state has $x=c:=2k_1/k_5$. On this plane $\dot y=0$ gives
$y=(k_1z+k_6c^2)/(k_4z)$, and then
\[
k_4z\,\dot z=-k_2k_4z^3+k_1k_4z^2+(k_1k_3-k_4k_6c^2)\,z+k_3k_6c^2 ,
\]
which is positive at $z=0$ and negative for large $z$. So a positive steady state exists for every $\kappa$, and ACR
holds for every $\kappa$. For the given rate constants $c=1$, $y=(5z+7)/(4z)$ and
$4z\,\dot z=-(2z-1)(z-1)(2z-7)$, which gives exactly $P_1,P_2,P_3$. The Jacobian matrices and their characteristic
polynomials are
\begin{align*}
J(P_1)&=\begin{psmallmatrix}-33&4&28\\14&-2&-14\\5&-1&-10\end{psmallmatrix}, &
\chi_1&=\lambda^3+45\lambda^2+206\lambda+30=(\lambda+5)(\lambda^2+40\lambda+6),\\
J(P_2)&=\begin{psmallmatrix}-38&8&14\\14&-4&-7\\10&-3&-4\end{psmallmatrix}, &
\chi_2&=\lambda^3+46\lambda^2+47\lambda-50,\\
J(P_3)&=\begin{psmallmatrix}-63&28&4\\14&-14&-2\\35&-13&-4\end{psmallmatrix}, &
\chi_3&=\lambda^3+81\lambda^2+632\lambda+1050 .
\end{align*}
By the Routh--Hurwitz criterion ($45\cdot206>30$ and $81\cdot632>1050$, all coefficients positive) $P_1$ and $P_3$
are exponentially stable; $\chi_2(0)=-50<0$ gives $P_2$ a positive real eigenvalue. All three steady states are
nondegenerate, so they persist, with their stability types, for nearby rate constants, and ACR holds for all rate
constants.
\end{proof}

The mechanism is a weighted integral feedback: the pair of complexes $\{Z,\,X+Z\}$ is the only one whose reactions
change $w\cdot x=x+2y$, so that $w\cdot f=z\,(2k_1-k_5x)$. Two further certified networks with six complexes, one of
them weakly reversible with two linkage classes of deficiency zero and total deficiency one, are given in the
sub-study \texttt{complexes/report.md}.

\section{Five reactant complexes suffice}\label{sec:five}
Consider four species $A,B,C,R$ and the network $\mathcal N_2$ of six reactions
\begin{equation}\label{eq:network}
\begin{gathered}
 0\xrightarrow{64}C,\qquad A+B\xrightarrow{3}2B,\qquad
 A+C\xrightarrow{5}B+C,\qquad B+C\xrightarrow{64}A,\\
 0\xrightarrow{1}R,\qquad R\xrightarrow{1}0.
\end{gathered}
\end{equation}
Its five distinct reactant complexes are $0,A+B,A+C,B+C,R$. Every reactant and product has
molecularity at most two. With the mass-action convention $k x^y$ its equations are
\begin{equation}\label{eq:exampleode}
\begin{aligned}
 \dot a&=-3ab-5ac+64bc, &\quad \dot b&=3ab+5ac-64bc,\\
 \dot c&=64-64bc, &\quad \dot r&=1-r.
\end{aligned}
\end{equation}
The stoichiometric subspace is $\{v:v_a+v_b=0\}$, of dimension three; the only independent
conservation law is $a+b=T$. Every positive equilibrium has $r=1$, so ACR holds globally, across
all compatibility classes. In fact arbitrary positive rates in the last two reactions give the
ACR value $k_{0R}/k_{R0}$.

In the class $a+b=9$, equilibrium requires $b=9-a$, $c=1/(9-a)$ and $r=1$, with $0<a<9$.
The remaining equation factors exactly as
\[
 \dot a\big|_{b=9-a,\,c=1/(9-a)}
 =\frac{-3(a-4)(a-6)(a-8)}{9-a}=0.
\]
Thus the complete list of positive equilibria in this class is
\[
 P_1=(4,5,1/5,1),\qquad P_2=(6,3,1/3,1),\qquad P_3=(8,1,1,1).
\]
In the class coordinates $(a,c,r)$ the restricted Jacobian is
\[
 J_{\rm red}=\begin{pmatrix}
 6a-27-69c&576-69a&0\\
 64c&-64(9-a)&0\\
 0&0&-1
 \end{pmatrix}.
\]
Its characteristic polynomial is $(\lambda+1)(\lambda^2-\tau\lambda+\Delta)$, where
$\tau$ and $\Delta$ are the trace and determinant of the upper left $2\times2$ block:
\begin{center}
\begin{tabular}{crrl}
\toprule
Equilibrium & $\tau$ & $\Delta$ & Stability in the class\\
\midrule
$P_1$ & $-1684/5$ & $1536$ & Exponentially stable\\
$P_2$ & $-206$ & $-768$ & Saddle\\
$P_3$ & $-112$ & $1536$ & Exponentially stable\\
\bottomrule
\end{tabular}
\end{center}
The trace-negative, determinant-positive criterion for a real $2\times2$ matrix proves the two
stability assertions exactly; no numerical eigenvalue calculation is used.

The regulator can also affect the other species without increasing the number of sources. Replace
$0\xrightarrow{64}C$, $0\xrightarrow1 R$, $R\xrightarrow1 0$ by
$0\xrightarrow{62}C$, $0\xrightarrow1 C+R$, $R\xrightarrow1 C$.
Only $\dot c$ changes, by adding $r-1$. The equilibria are unchanged and the Jacobian remains
triangular in the two blocks, with exactly the same characteristic polynomial. Both versions are
verified directly from their reaction lists by \texttt{certify\_five\_example.py}.

\section{Four reactant complexes do not suffice}\label{sec:four}
\begin{theorem}\label{thm:rank-three}
Let a mass-action system have stoichiometric rank three, four distinct reactant complexes, and
at-most-bimolecular sources and products. If it has ACR, each positive compatibility class contains at most
two nondegenerate positive equilibria. When there are two, their restricted Jacobian determinants have
opposite signs. In particular, at most one is exponentially stable.
\end{theorem}

\begin{corollary}\label{cor:five}
An at-most-bimolecular mass-action system with ACR and two exponentially stable positive equilibria in one
class has at least five distinct reactant complexes.
\end{corollary}
\begin{proof}
With $r$ reactant complexes a positive equilibrium gives a strictly positive vector in the kernel of
the aggregated coefficient matrix $N$, so $\rank N\le r-1$. At a nondegenerate equilibrium,
$\im N$ equals the stoichiometric subspace $S$: indeed $\im J\subseteq\im N\subseteq S$, while
$J|_S$ is invertible. For $r\le4$ this forces $\dim S\le3$. Dimension two is excluded by
\zcref{thm:dim-two}, which holds for every number of species and reactant
complexes. Dimension one is excluded by Proposition~3.10 of~\autocite{KNS2024}; directly, on a class the
field is a fixed vector times a polynomial of degree at most two in the class coordinate, and two
simple zeros of such a polynomial have opposite derivative signs. Dimension zero has singleton classes.
Dimension three forces $r=4$ and is excluded by \zcref{thm:rank-three}.
\end{proof}
\zcref{thm:rank-three} does not use \zcref{thm:dim-two}; only \zcref{cor:five} does.

\section{Toric reduction and the determinant sign}
List the sources as $y_1,\ldots,y_4$, and write $f(x)=N m(x)$, where $m_a(x)=x^{y_a}$.
It suffices to consider a system possessing a nondegenerate positive equilibrium $x^*$.
Then $\rank N=3$ and $\ker N$ is generated by $m(x^*)>0$.
Let $A$ be the $3\times n$ matrix with rows $y_a-y_4$, $a=1,2,3$.
The equilibrium equations are
\[
 A\log x=A\log x^*,\qquad
 \mathcal E=x^*\circ\exp(\ker A).
\]
Nondegeneracy forces $\rank A=3$. Otherwise the differential of these equations, and hence
the restricted Jacobian, has rank at most two. ACR in species $i$ is equivalent to
$e_i\in\operatorname{row}(A)$, or to a zero $i$th row in any basis matrix $V$ of $\ker A$.

Let $W$ have rows forming a basis of $S^\perp$. Compatibility is $Wx=Wx^*$.
Set $N_0=[N_1\ N_2\ N_3]$ and $p_a=m_a(x^*)/m_4(x^*)>0$.
Any three columns of $N$ are independent. At equilibrium,
\[
 J=m_4(x)N_0\diag(p_1,p_2,p_3)A\diag(1/x).
\]
Consequently, for a fixed nonzero constant $\sigma$ and a positive factor $q(x)$,
\begin{equation}\label{eq:orientation}
 \det(J|_S)=\sigma q(x)
 \det\begin{pmatrix}W\\ A\diag(1/x)\end{pmatrix}.
\end{equation}
To see this, use the basis $N_0$ for $S$ and complete it by columns $B$ to a basis of $\R^n$.
Multiplication of the displayed block matrix by $[B\ N_0]$ gives a block triangular matrix with
diagonal blocks $WB$ and $A\diag(1/x)N_0$. All change-of-basis determinants are fixed and nonzero.

The same determinant can be written with a basis matrix $V$ of $\ker A$. Choose $C$ with $AC$ invertible.
Then $[V\ C]$ is invertible, and multiplying the block matrix on the right by $\diag(x)[V\ C]$ gives a block
triangular matrix with diagonal blocks $W\diag(x)V$ and $AC$. Hence, for a fixed constant $c\ne0$,
\begin{equation}\label{eq:orientationV}
 \det\begin{pmatrix}W\\ A\diag(1/x)\end{pmatrix}
 =\frac{c}{x_1\cdots x_n}\,\det\bigl(W\diag(x)V\bigr).
\end{equation}
At a positive equilibrium, nondegeneracy is therefore equivalent to $\det(W\diag(x)V)\ne0$, and the sign
of $\det(J|_S)$ is a fixed sign times the sign of $\det(W\diag(x)V)$.

Replacing logarithmic binomials by ordinary binomials multiplies their differential rows by positive
monomial factors, followed by a fixed nonsingular row transformation. Fixing the ACR concentration and
performing linear eliminations therefore preserves the comparison of determinant signs.
In a one-variable reduction $F(t)=0$, the augmented determinant is a fixed sign times a positive
factor times $F'(t)$, provided the eliminated equations have nonsingular differentials on the positive
parameter interval. We verify this condition explicitly below. Two simple roots of a quadratic on an
interval have opposite derivative signs. The same holds for $h(t)-p$ when $p>0$ and $\log h$ is
strictly concave on that interval. Since a real Hurwitz $3\times3$ matrix has negative determinant,
opposite determinant signs preclude two stable equilibria.

\section{Exhaustive source classification}
A species absent from every source has a nonnegative derivative. At a positive equilibrium every
reaction rate is positive, so no reaction can produce that species either. Such species are inert and
can be removed. They cannot be ACR species, since their arbitrary positive concentrations give
equilibria with different values. Four bimolecular sources then involve at most eight species.

The source classification is elementary and finite. For $4\le n\le8$, enumerate four-element subsets
of
\[
 \{0,e_i,2e_i,e_i+e_j:1\le i<j\le n\}.
\]
Retain only sets covering all species, having $\rank A=3$, and having an ACR coordinate.
The rank test uses integer $3\times3$ minors. The ACR coordinates are precisely the indices contained
in every nonzero three-column minor: deleting such a column lowers the rank to two.
Canonicalization permutes the four sources and sorts the species columns, so no floating-point
rank decision or heuristic is involved. The script \texttt{certify\_sources.py} performs this enumeration.

For three species the toric equilibrium set is a singleton. For four species there are 26 source types.
In each, a generator $v$ of $\ker A$ has at most three distinct entries, including zero; this is checked
by its integer cofactor formula. The class equation is
\[
 \sum_{j=1}^4 W_j x_j^* e^{v_j t}-Wx^*=0.
\]
It has at most three distinct exponential terms. A nonzero sum of three real exponentials has at most
two real zeros counted with multiplicity: divide by the smallest exponential and apply Rolle's theorem
to the resulting derivative, which has at most two terms. If the sum is identically zero, all equilibria
on the curve are degenerate. In the nondegenerate case the two determinants have opposite signs by
\eqref{eq:orientation}.

For five or more species the exact list, up to species and source permutations, is as follows.
\begin{center}
\begin{tabular}{cll}
\toprule
Type & Sources & ACR species\\
\midrule
S1 & $0,E,C+D,A+B$ & $E$\\
S2 & $0,2E,C+D,A+B$ & $E$\\
S3 & $E,2E,C+D,A+B$ & $E$\\
S4 & $E,D,C+E,A+B$ & $C$\\
S5 & $E,D+E,C+E,A+B$ & $C,D$\\
S6 & $E,D+E,C+D,A+B$ & $D$\\
S7 & $E,D+E,2C,A+B$ & $D$\\
S8 & $E,D+E,B+C,A+C$ & $D$\\
S9 & $F,E+F,C+D,A+B$ & $E$\\
\bottomrule
\end{tabular}
\end{center}
There are eight types on five species, one on six species, and none on seven or eight.

For S4, S5, S6 and S8 the varying equilibrium concentrations, after positive constant factors are
removed, belong to three groups $g_1=u,g_2=w/u,g_3=w$. For S4, $c$ is fixed and $d,e$ are
proportional to $ab$; for S5, $c,d$ are fixed and $e$ is proportional to $ab$; for S6, $d$ is fixed
and $c,e$ are proportional to $ab$; for S8, $d$ is fixed, $a,b$ are proportional, and $e$ is
proportional to $ac$. The class equations are $Qg=T$, two linear equations in the three group
variables. If $\rank Q<2$, the equilibria are degenerate. Otherwise write the affine solution line
as $g(t)=g^0+t\rho$, where $\rho$ spans $\ker Q$. Positivity of all three components cuts out an
interval. The remaining equation is the polynomial $F(t)=g_1(t)g_2(t)-g_3(t)$ of degree at most two.
The augmented determinant with rows $Q$ and $\nabla(g_1g_2-g_3)$ is a fixed nonzero multiple of
$F'(t)$: the cross product of the rows of $Q$ is a fixed nonzero multiple of $\rho$.
The group embedding is linear with independent columns and positive constant entries on disjoint
coordinate groups. Eliminating the other fixed linear equations therefore changes the original
augmented determinant only by a fixed nonzero factor. Hence \eqref{eq:orientation} applies, and
two nondegenerate equilibria have opposite determinant signs. If $F$ vanishes identically, every
equilibrium is degenerate.

\section{Exhaustive stoichiometric classification}
We now give an exact classification for S1, S2, S3, S7 and S9.

For a fixed source set, form every possible nonzero reaction vector $p-y_a$, where $p$ is any
at-most-bimolecular product. Every rank-three stoichiometric subspace is spanned by three of these
vectors. Enumerate these triples and identify equal spans by their primitive exterior coordinates
(their $3\times3$ minors, divided by their gcd, with the first nonzero entry positive).
For each span $S$, retain all allowed reaction vectors lying in it, giving a maximal matrix $G$.
Membership is checked by the integer adjugate identity in a nonsingular three-coordinate projection.

\begin{lemma}\label{lem:maximal}
If some subnetwork with this span has a positive equilibrium, the maximal matrix $G$ has a strictly
positive kernel vector. In a three-coordinate projection, failure of this condition is detected by
a nonzero cross product of two columns having nonnegative inner products with every column,
after possibly reversing its sign.
\end{lemma}
\begin{proof}
The subnetwork equilibrium supplies a positive dependence among vectors spanning $S$. Their positive
cone is therefore all of $S$: any linear representation can be made nonnegative by adding a sufficiently
large multiple of that dependence. Each additional vector can thus be cancelled by a nonnegative
combination of the old vectors. Adding these dependencies to the original positive dependence gives
a dependence positive on every maximal-network reaction.

Conversely, a finite collection spanning $\R^3$ positively spans $\R^3$ exactly when its dual cone is
$\{0\}$. If this dual cone is nonzero it is pointed, since the original vectors span $\R^3$.
It has an extreme ray, on which at least two independent constraints are active. Such a ray is
perpendicular to two independent original vectors, and hence is their cross product up to sign.
This proves the exact finite test.
\end{proof}

Next write
\begin{equation}\label{eq:minors}
 D_W(z)=\det(W\diag(z)V)
       =\sum_{|I|=n-3}\det W_{:,I}\det V_{I,:}\prod_{i\in I}z_i.
\end{equation}
If all coefficients vanish, every equilibrium is degenerate by \eqref{eq:orientation} and
\eqref{eq:orientationV}. If the nonzero coefficients have one
sign, there is at most one positive equilibrium per class. Indeed, if $x,y$ are two such equilibria,
$\log x-\log y=Vt$ and $x-y=\diag(\ell)Vt$ with all logarithmic means $\ell_i>0$.
Then $W\diag(\ell)Vt=0$, contradicting \eqref{eq:minors} unless $t=0$.

A further exact exclusion concerns the aggregated kinetic rank. Let $K$ be a rational basis matrix
of $\ker G$, write a general flux as $K\theta$, and sum its reaction vectors separately for each
source, giving columns $L_1(\theta),\ldots,L_4(\theta)$ with sum zero. Project onto three independent
coordinates. If the cubic polynomial
\[
 \det[L_1(\theta)\ L_2(\theta)\ L_3(\theta)]
\]
vanishes identically, $\rank N<3$ at every positive equilibrium of every subnetwork with this span.
Such equilibria are degenerate. The certificate checks polynomial identity exactly. Nonzero evaluations
are used only to retain a case, never to discard one.

The complete counts are:
\begin{center}\small
\begin{tabular}{crrrrr}
\toprule
 & S1 & S2 & S3 & S7 & S9\\
\midrule
All rank-three spans & 4333&4287&4333&4397&23035\\
Missing a source &2965&3161&2965&3050&17734\\
No positive maximal flux &1126&782&1126&1042&4442\\
Injective &169&301&129&231&569\\
Class-degenerate &34&6&6&32&144\\
Kinetically degenerate &39&37&43&20&92\\
Remaining &0&0&64&22&54\\
\bottomrule
\end{tabular}
\end{center}
Each span occurs in exactly one row after the first row. After fixing the ACR coordinate, remove its
column from $W$ and row-reduce. Permutations $A\leftrightarrow B$, $C\leftrightarrow D$, and interchange
of the two pairs are allowed for S3 and S9; only $A\leftrightarrow B$ is needed for S7. The remaining
spaces reduce to eleven, five and three forms respectively. The next sections list \emph{all} of them
and prove the required root and determinant properties.

\section{The eleven S3 forms}
The positive equilibria have
\[
 e=e_0,\qquad ab=p,\qquad cd=q,\qquad p,q,e_0>0.
\]
After $e$ is fixed the two conservation equations have one of the following forms. Constants $T,U$
are arbitrary real numbers; existence of a positive solution supplies any needed positivity restrictions.
\begin{center}\small
\begin{tabular}{cll}
\toprule
Index & First equation & Second equation\\
\midrule
0 &$b=T$ &$c+d=U$\\
1 &$a-2c+2d=T$ &$b+c=U$\\
2 &$a-c+d=T$ &$b+(c+d)/2=U$\\
3 &$a-c+d=T$ &$b+c=U$\\
4 &$a-c+d=T$ &$b+c+d=U$\\
5 &$a-c+d=T$ &$b+2c=U$\\
6 &$a-c/2+d=T$ &$b+c/2=U$\\
7 &$a=T$ &$b+c+d=U$\\
8 &$a+d=T$ &$b+c=U$\\
9 &$a+(c+d)/2=T$ &$b+(c+d)/2=U$\\
10 &$a+c+d=T$ &$b+c+d=U$\\
\bottomrule
\end{tabular}
\end{center}
Forms 0 and 7 fix $a,b>0$ and $c+d$, so $cd=q$ is a quadratic on the interval $c,d>0$.
For form 8 take $a=t$, $b=p/t$, $c=U-p/t$, $d=T-t$. The positive interval is
$p/U<t<T$, with $T,U>0$. Multiplication of $cd-q=0$ by $t>0$ gives
\[
 -Ut^2+(UT+p-q)t-pT=0.
\]
For forms 9 and 10, subtraction fixes $a-b$; together with $ab=p$ this uniquely fixes the positive
pair $(a,b)$, since $a-p/a$ is strictly increasing from $-\infty$ to $+\infty$.
Again $c+d$ is fixed and $cd=q$ is quadratic on its positive interval.

For forms 1--6 put $c=t>0$, $d=q/t$. Then
\[
 a=T+\alpha t-\beta q/t,\qquad
 b=U-\gamma t-\delta q/t,
\]
where $\alpha,\beta,\gamma>0$ and $\delta\ge0$, as read directly from the table.
The functions $a,b$ are concave; $a$ is strictly concave. Their common positivity set is an interval.
On it,
\[
 (\log(ab))''=
 \frac{a''}{a}-\frac{(a')^2}{a^2}
 +\frac{b''}{b}-\frac{(b')^2}{b^2}<0.
\]
Hence $ab=p$ has at most two roots, with opposite derivative signs when two roots exist.

For completeness, take the augmented determinant with rows $R$, $\nabla(ab-p)$, $\nabla(cd-q)$,
where $R$ is the conservation matrix in the table. On the displayed scalar curves its multiplier of
the derivative is, up to a fixed sign: $T>0$ for forms 0 and 7; $t>0$ for forms 1--6 and 8;
and $h(a+b)>0$ for forms 9 and 10, where $h=1/2$ and $h=1$ respectively.
For form 8 multiplication of the scalar equation by $t$ preserves the sign at a root.
These identities also prove that repeated roots and identically zero reductions cannot be
nondegenerate equilibria.

\section{The five S7 forms}
Here the positive equilibrium equations are
\[
 d=d_0,\qquad ab=pe,\qquad c^2=qe,\qquad p,q,d_0>0.
\]
The reduced variables are $(a,b,c,e)$, and the conservation equations are:
\begin{center}\small
\begin{tabular}{cll}
\toprule
Index & First equation & Second equation\\
\midrule
0 &$b+e=T$ &$c-e=U$\\
1 &$a-c+2e=T$ &$b+c=U$\\
2 &$a-c/2+e/2=T$ &$b+c/2+e/2=U$\\
3 &$a-e=T$ &$b=U$\\
4 &$a+b+e=T$ &$c=U$\\
\bottomrule
\end{tabular}
\end{center}
For form 0 take $c=t$, $e=t-U$, $b=T-t+U$, $a=pe/b$.
Positivity is an interval, and $c^2-qe=t^2-q(t-U)=0$ is quadratic.
Its determinant multiplier has magnitude $b>0$.

For forms 1 and 2 put $u=1/c>0$, so $e=1/(qu^2)$. Then
\[
 A(u):=a/c=Tu+\alpha-\frac{\beta}{qu},\qquad
 B(u):=b/c=Uu-\gamma-\frac{\delta}{qu},
\]
with $(\alpha,\beta,\gamma,\delta)=(1,2,1,0)$ or $(1/2,1/2,1/2,1/2)$.
The positive interval is $A(u)>0$, $B(u)>0$. Both functions are concave and $A''(u)<0$.
Thus $\log(A B)$ is strictly concave, and $A B=p/q$ has at most two roots with opposite
derivative signs. If $H=AB-p/q$, the augmented determinant with rows $R$, $\nabla(ab-pe)$,
$\nabla(c^2-qe)$ equals, up to fixed orientation,
\[
 q\bigl(H'(u)-2H(u)/u\bigr).
\]
At a root its multiplier of $H'$ is the positive constant $q$.

For form 3 the equation $ab=pe$ is $UT+(U-p)e=0$, with $U=b>0$.
There is at most one positive solution, unless the equation is an identity, in which case the
positive solutions form a curve and are degenerate. For form 4, $c=U>0$ fixes $e=U^2/q$;
the sum and product of $a,b$ are fixed, giving a quadratic on $a,b>0$.
The determinant multiplier has magnitude $q>0$ in both scalar parametrizations used by the certificate.

\section{The three S9 forms}
The equilibrium equations are
\[
 e=e_0,\qquad ab=pf,\qquad cd=qf,\qquad p,q,e_0>0.
\]
In variables $(a,b,c,d,f)$ the three conservation equations have one of these forms:
\begin{center}\small
\begin{tabular}{clll}
\toprule
Index & First equation & Second equation & Third equation\\
\midrule
0 &$b+f=T$ &$c-f=U$ &$d=V$\\
1 &$a-c/2+f/2=T$ &$b+c/2+f/2=U$ &$d=V$\\
2 &$a=T$ &$b=U$ &$c+d+f=V$\\
\bottomrule
\end{tabular}
\end{center}
For form 0, put $f=t$, $b=T-t$, $c=t+U$, $d=V$, $a=pt/(T-t)$.
Positivity is an interval with $V>0$. The remaining equation is
$(V-q)t+VU=0$. It has at most one solution, or is an identity giving only degenerate equilibria.
The determinant multiplier has magnitude $b>0$.

For form 1, put $f=t$, $d=V>0$, $c=qt/V$,
\[
 a=T+\tfrac12(q/V-1)t,\qquad
 b=U-\tfrac12(q/V+1)t.
\]
The positive domain is an interval defined by three linear inequalities $t,a,b>0$.
The remaining equation $ab-pt=0$ is quadratic. Its determinant multiplier has magnitude $V>0$.

For form 2, $a=T>0$ and $b=U>0$ fix $f=TU/p$.
The sum $c+d=V-TU/p$ and product $cd=qTU/p$ are fixed, so there are at most two positive
solutions, with opposite derivative signs. The determinant multiplier has magnitude $p>0$.

All nineteen forms have now been treated. Together with the source classification, the injectivity
and degeneracy exclusions, and \eqref{eq:orientation}, this proves \zcref{thm:rank-three}.

\section{Further results and open questions}\label{sec:further}

The following results complement \zcref{tab:summary}. Their proofs, partly computer-assisted with exact arithmetic, are
in the two sub-study reports \texttt{examples/report.md} and \texttt{complexes/report.md}, with the scripts that
check them.

\begin{enumerate}
\item \emph{A second three-species example with rational steady states.} The network
$\mathcal G_7=\{X+Y\xrightarrow{1}2Y,\ Y\xrightarrow{1}2Z,\ 0\xrightarrow{5}X+Z,\ X+Z\xrightarrow{4}2Z,\
2Z\xrightarrow{1}X+Z,\ Y+Z\xrightarrow{6}Y\}$ has $\dot y=y(k_1x-k_2)$, hence ACR in $X$ for all rate constants,
and for the indicated ones exactly three positive steady states $(1,2,1)$, $(1,\frac54,\frac32)$,
$(1,\frac{10}9,\frac53)$, the first and the last exponentially stable (exact Gröbner basis and Routh--Hurwitz).
It has $7$ complexes.
\item \emph{Six complexes in two linkage classes.} The weakly reversible network
$\{2Z\rightleftarrows X+Z,\ Y+Z\to X,\ Y+Z\to X+Y,\ X\to2X,\ X+Y\to X,\ 2X\to Y+Z\}$ has $6$ complexes, two linkage
classes of deficiency~$0$ and total deficiency~$1$, ACR in $X$ and, for suitable rate constants, two exponentially
stable positive steady states. So the deficiency-one theorem cannot be applied to two linkage classes when
$\sum_\theta\delta_\theta<\delta$, and six-complex examples need not have a single linkage class.
\item \emph{Reactions with three species.} No full-dimensional bimolecular network with $3$ species and at most
$5$ reactions has three nondegenerate positive steady states, or two exponentially stable ones, for any rate
constants. Networks with at most $4$ distinct reactant complexes are covered by
\autocite[Propositions~3.1(3) and~3.5]{KNS2024}; the $3\,188\,646$ networks with $5$ distinct reactant complexes
were enumerated exactly, through a reduction to the sign changes of one univariate function whose derivative has,
up to a fixed sign, the sign of the Jacobian determinant (two exponentially stable states have determinants of the same sign, which
forces a third steady state between them). By \zcref{cor:dimension}, a three-species network with ACR and
multistability is full-dimensional, so it needs at least $6$ reactions, which $\mathcal N_1$ and $\mathcal G_7$
attain. The restriction to full-dimensional networks is necessary: the two-dimensional network
$\{0\xrightarrow{10}Z,\ Y+Z\xrightarrow{10}X,\ X+Z\xrightarrow{1}Y+Z,\ X+Y\xrightarrow{1}2Y\}$ has the three
nondegenerate positive steady states $(5,1,1)$, $(4,2,\frac12)$, $(3,3,\frac13)$ in the class $x+y=6$, the outer two
exponentially stable (it has no ACR).
\end{enumerate}

\paragraph{Open questions.} Does one network attain all four minima of \zcref{tab:summary}? $\mathcal N_1$ has six
reactant complexes and $\mathcal N_2$ has four species. In particular, for full-dimensional networks on three
species, is the bound $n+2=5$ of \autocite[Theorem~5.1]{KNS2024} attained with multistability? Does
\zcref{thm:dim-two} hold for networks with bimolecular reactant complexes but larger products
(\zcref{rem:products})? Finally, \autocite{JNJ2026} found ACR with bistability in a dual-site futile cycle with
bifunctional enzymes (six species, dimension four); our examples show that much smaller mechanisms, based on a
weighted integral feedback, suffice.

\section{Computer verification}\label{sec:computer}

All scripts print \texttt{PASS}/\texttt{FAIL} for each claim and exit with a nonzero code on failure. They were run
under WSL (Fedora~44) with Python~3.12, sympy~1.14, python-flint~0.9 and NumPy.
\begin{itemize}
\item \texttt{scripts/check\_dimension\_two\_example.py} (8~s): on $20$ random rational instances of the KNS network
of Example~4.10, checks the determinant formula of \zcref{lem:det} and that the two positive steady states of a
class have determinants of opposite signs. This is a sanity check; \zcref{thm:dim-two} is proved by hand.
\item \texttt{complexes/scripts/certify\_examples.py} (exact rational arithmetic, standard library only, under
$0.1$~s): \zcref{thm:three-species} and the two further six-complex examples (ACR identity, steady states,
characteristic polynomials, Routh--Hurwitz conditions, counts of complexes and deficiency).
\texttt{complexes/scripts/eigen\_enclosures.py}: eigenvalue enclosures in ball arithmetic.
\item \texttt{examples/scripts/certify\_example.py} (under $1$~s): the network $\mathcal G_7$ (exact Gröbner basis,
root isolation, ball-arithmetic eigenvalues). The enumeration of item~3 is
\texttt{examples/scripts/enum5\_reactions.py} (about two minutes on $12$ cores).
\end{itemize}
Searches and simulations used to find the examples are in \texttt{explore} sub-folders; they are floating-point and
not part of any proof.

For the four-reactant classification and the network $\mathcal N_2$, run from the folder \texttt{four-reactants/}:
\begin{verbatim}
python scripts/certify_sources.py
python scripts/certify_subspaces.py
python scripts/certify_normal_forms.py
python scripts/certify_five_example.py
\end{verbatim}
The first script re-enumerates every eligible source set. The second generates from scratch all
40,385 reaction spans for the five specified source types, proves each exclusion, and writes the
retained spaces and their classifications to \texttt{certificates/subspaces.json}. It does not read
any numerical candidate list. The third symbolically checks all nineteen determinant identities
and the concavity derivatives, including the conservation equations along every scalar curve.
The fourth derives the upper-bound example from its reaction list, factors its complete steady-state
equation and checks every stability claim in rational arithmetic.
Each script prints PASS/FAIL and exits nonzero on a failure.

The environment used was Python 3.12, SymPy 1.14 and NumPy under Fedora 44 in WSL. The runs take
less than one minute in total. Integer arrays use signed 64-bit entries; their input values are
in $[-2,2]$, three-minors have absolute value at most 48, adjugate entries at most 8, membership
products at most 288, and cross-product inner products at most 48. Thus these operations cannot
overflow. Nullspaces, polynomial identities and all subsequent coefficients are rational and symbolic.
No numerical eigenvalue, optimizer tolerance, random search or sampling result is used as an exclusion.

Fresh-context verifiers checked the two reports merged here; see \texttt{verification.md} and
\texttt{four-reactants/verification.md}. The second one reimplemented the enumeration, obtained the same label for
every span, and recorded two facts that make the proof sturdier. The positive-flux exclusion is not needed: every
span it rejects is also rejected by the injectivity or degeneracy tests. The kinetic exclusion also has an elementary
reason, without the cubic: on each span that it removes, the spans of the reaction vectors of some three sources
fail Rado's condition, so three columns of $N$ are forced to be dependent.

The merged text itself has not been re-verified. A public-source novelty audit on 31 October 2026, recorded in
\texttt{four-reactants/references.md}, found no prior resolution of the exact-minimum question; absence of search hits
is not a proof of priority. This is a computer-assisted proof, checked by fresh-context AI verifiers, not by human
referees. The results do not concern larger products or nonhyperbolic notions of stability, and they do not classify
every network attaining the minima.

\paperreferences
\end{document}
